\MajorSection{الخاتمة}

هنا يختم الحاج دراسته العملية للبشر. وصورة القدر الذي يلعب بالناس قطعًا من اللعب نظرة شائعة عند الشرقيين. وفكرته عن الحكمة هي أيضًا فكرة بوب:

\begin{SourceQuotation}
وكل معرفتنا أن نعرف أنفسنا.\\
المقالة الرابعة، البيت ٣٩٨.
\end{SourceQuotation}

كان الأسف، أي التوبة، واحدة من الخطايا المهلكة الاثنتين والأربعين عند المصريين القدماء. وتقول شعائر الموتى، في تبرئة النفس أو الشبح بالنفي: «لا تأكل قلبك»، كما في كتاب لبسيوس «أقدم نصوص كتاب الموتى». وقد استعرنا الامتحان التنافسي من الصينيين؛ وفي أيامنا المرضية هذه من التأمل الضعيف إلى الداخل والوراء، ربما تعلّمنا الحكمة من المصريين القدماء الأشداء. وحين يغني «انصرف عن لماذا واطلب كيف»، يشير إلى التمييز المدرسي القديم بين البرهان على العلة، \textenglish{Demonstratio propter quid}، أي لماذا يكون الشيء، وبين البرهان على وجوده، \textenglish{Demonstratio quia}، أي أن الشيء موجود. وينتهي «الإنسان العظيم» إلى أن يصير بلا موت، كما يقول شكسبير في سونيتته النبيلة:

\begin{SourceQuotation}
وحين يموت الموت، لن يكون ثمة موت بعد ذلك!
\end{SourceQuotation}

يرى الحاج، مثل الوثنيين العظماء، أن الإنسان ولد خيّرًا؛ بينما يتمسك المسيحي، «المعذّب بالأمور الإلهية»، بالعقيدة المعزّية في الخطيئة الفطرية. ومن هنا المبدأ العام أن على الإنسان فعل الخير ليكسب به هنا أو في الآخرة؛ تلك «الأنانية المستنيرة» التي تقول: أحسن الفعل، واحصل على فائدة مركبة في حال مستقبلية. ويبدو أن الإشارة إلى «كلمة المؤمن بإله شخصي» تعني أن أتباع إله شخصي يجب أن يؤمنوا بالعلم السابق المطلق لمحيط العلم، في الجزئيات كما في الكليات. وسيادة الناموس تحرر الإنسان؛ واستثناءاته هي الفجوات التي يتركها جهله. والنواح على الزهرة الساقطة وغيرها يذكّرنا بالبُلامبال، المراثي، للكاتب المناهض للبرهمانية باتيرا گيريار. والإشارة إلى مايا مأخوذة من داس كبير:

\begin{SourceQuotation}
تموت مايا، ولا يموت العقل؛ وقد مات الجسد ومضى.\\
\textenglish{Mâyâ mare, na man mare, mar mar gayâ, sarîr}.
\end{SourceQuotation}

لقد قلت إن النيرفانا انطفاء جزئي بالاندماج في الأعلى، ولا ينبغي خلطها بالبارينيرفانا، أو الفناء المطلق. وفي الأولى أيضًا يلد الموت كائنًا جديدًا، هو تجسيد الكارما، أي الأعمال الخيّرة والشريرة التي أُنجزت في عصور التناسخ التي لا تُحصى.

هنا ينتهي نصيبي من العمل. وقد كان، في مجموعه، كبيرًا. حذفت، كما رأينا، مقاطع مختلفة، وغيّرت ترتيب مقاطع أخرى. ولم يُترجم النص حرفيًا في أي موضع؛ بل أُعطيت في الواقع صبغة أوروبية مألوفة لكثير من المشاعر التي عُدّت شرقية أكثر مما ينبغي. ولما كان الوزن الذي اختاره الحاج عبدو هو البحر الطويل، رأيت من المناسب حفظ هذه الخصوصية، وتزيين أطرافه بقافية الأصل الخشنة غير اللافتة.

عش، وكن بخير!
